\documentclass{beamer}

\mode<presentation> {

\usetheme{Szeged}
\usecolortheme{whale}
\usepackage{url}
\usepackage{xcolor}
\usepackage{multicol}
\setbeamertemplate{footline}[page number] 
}

\usepackage{graphicx} 
\usepackage{booktabs}
\usepackage{amsmath}
\usepackage{subfig}
\usepackage{caption}
\usepackage{mathabx}
\usepackage{wrapfig}
\usepackage{tikz}
\usepackage{animate}
\usepackage{multicol}
\usetikzlibrary{shapes.geometric, arrows}
\usepackage{calligra}
\usepackage[T1]{fontenc}
\usepackage{algorithm}
\usepackage{algorithmic}

\usepackage{diagbox}

\newcommand{\LINEIF}[2]
    {%
        \STATE\algorithmicif\ {#1}\ \algorithmicthen\ {#2} \algorithmicendif
    }

\title{Inventory Routing Problem with Logistic Ratio}

\date{\today}



\begin{document}




\begin{frame}{Progress 12th Nov 2024}

\begin{itemize}
    \item Implemented cheapest insertion. I believe it is an efficient implementation through some data structure.
    \item Local search has been updated
    \item Update ways to adjust the penalty on stock out.
    \item With just one operator, it can find feasible solutions that improve the initial logistic ratio. 
    \item 1000 iterations on instance/mvirp/Small$\_$Instances/Istanze0105h6/$abs5n50\_4.dat$ took around 258 seconds.

    
    \item Take note of parameters tested. And tune them by \url{https://mlopez-ibanez.github.io/irace/}
    
 
\end{itemize}
    
\end{frame}


\begin{frame}{Progress 4th Dec 2024}

\begin{itemize}
    \item The swap operator can now obtain more feasible solutions 
    \item The efficiency in evaluating the solutions has been significantly improved.
    \item These improvements will also facilitate the implementation of other upcoming operators.
\end{itemize}
On instance Istanze0105h6/$abs5n50\_4.dat$
\begin{itemize}
    \item Running 1000 times, the solution is improved by $20\%$ of the initial solution. All it requires is 10 seconds. 
\end{itemize}
Upcoming operators
\begin{itemize}
    \item Extend the current swap operator by purely removing an allocated customer
    \item Operator that swaps between two allocated customers
    \item Increase a visit to a customer at the stockout period.
    \item Remove a visit
    \item Balancing quantity operator
\end{itemize}
\end{frame}

\begin{frame}{Progress 28th Jan 2025}

\begin{itemize}
    \item Take any schedule, the operator does the following things.
    \begin{itemize} \setlength{\itemindent}{1em}
        \item [Step 1]  Reset the delivery quantity and inventory level.
        \item [Step 2] Starting from the first time period.
        \item [Step 3] Randomly pick a customer from those who are visited.
        \item [Step 4] Deliver as much as possible to this customer with respect to the vehicle capacity and maximum inventory level.
        \item [Step 5] Update the inventory level.
        \item [Step 6] Go to \textbf{Step 3} until no customer to pick.
        \item [Step 7] Go to \textbf{Step 2} until the last time period.
    \end{itemize}
    \item Issues:
    \begin{itemize}
        \item The rebalance operator may give a solution with a stockout.
        \item Customers who got picked up later may not deliver anything due to vehicle capacity.
    \end{itemize}  
\end{itemize}
\end{frame}

\begin{frame}{Progress 11th Feb 2025}

The adaptive weighted probability is a matrix, a row represents a customer and a column represents a time periods. It is comprised of the following components:
\begin{itemize}
    \item Before executing any algorithm, run a check on the data and compute $\alpha\times\frac{Inventory\_Begin}{Inventory\_Max} + (1-\alpha) \frac{Inventory\_Max}{Demand}$
    
    $\frac{Inventory\_begin}{Demand(\text{Next visit time - current visit time})}$
    
    \item Takes into account the input schedule, if a customer has not been visited, the probability of selecting this customer becomes 0.
    \item  
    
    $Weight = \begin{cases}
        Weight\times 2, & \text{If not visited the next day} \\
        \frac{Weight}{2} , & \text{if visited the next day}
        \end{cases}$
    \item 
    $Weight = Weight \times (1 - \frac{\#\text{ of visits}}{\#\text{ of days }+1}) $
    
\end{itemize}
\end{frame}
\begin{frame}{Progress 25th Feb 2025}

Some behaviour observed
\begin{itemize}
    \item Apply a procedure once that delivering the minimum quantity across all periods is not working. As the demand accumulates to later periods, it leads to infeasibility. Changes were made accordingly.
    \item Some moves may result in an obvious infeasible solution, so obviously the rebalance will return an infeasible solution. There is a check, the rebalance skips the solution if it is true.
    \item Roughly speaking, within the neighbourhood of the swap operator, for 70$\%$ of the neighbour, the rebalance can make them feasible.
    \item After rebalance some visits can be removed without incurring stockout. In my opinion, it happens quite frequently.
    
\end{itemize}
\end{frame}

\begin{frame}{Progress}
\begin{itemize}
    \item The insert part of the operator could fail due to a full vehicle. Swap part works since the inserted customer delivers at least the quantity removed.
    \item Current implementation doesn't insert customer into empty route.
    \item The implementation for evaluating a solution still need improve.
\end{itemize}
\end{frame}


\begin{frame}{Recall}
\begin{itemize}
    \item Also implement swap between 2 routes.
\end{itemize}
\end{frame}

\begin{frame}{Progress}
\begin{itemize}
    \item Previously, the disturbances randomly remove some routes. This procedure does not produce enough variation. In the updated disturbances, it randomly removes a sequence of consecutive customers. The length of the sequence is also randomly chosen.
\end{itemize}
\end{frame}
\end{document}